% ===================================================================
%  BẢNG Số 8 — Mùa gặt.
%
%  Traduction, faite en 2026, en vietnamien standard d'aujourd'hui,
%  sur l'ido de texto/io. Regles de la colonne : voir 10-bang-01.tex.
%  Deux scenes, deux numerotations : les cles portent c1 et c2.
%
%  LA NOTE (*) DU LIVRET RECOIT ICI UNE REPONSE QU'ELLE N'ATTENDAIT
%  PAS. Guignon s'y plaint que « rekolto » soit imprecis — recolte du
%  lin, du raisin, des pommes ? — et propose une racine nouvelle pour
%  la seule moisson des cereales. Or le vietnamien fait exactement ce
%  qu'il demande, et depuis toujours : « gặt » ne se dit QUE de la
%  cereale coupee a la faucille, « hái » de ce qu'on cueille a la
%  main, « thu hoạch » de la recolte en general. La langue d'un pays
%  de riz avait separe ce que l'ido cherchait a separer.
%
%  ET LE VIGNOBLE, LUI, N'A PAS DE MOTS. La seconde scene est celle
%  des vendanges, et c'est le passage du livret le plus etranger a ce
%  pays : ni la vigne, ni le pressoir, ni le tonnelier, ni le cep
%  n'ont de nom ancien en vietnamien. On les rend par des composes
%  transparents — « máy ép nho » la machine a presser le raisin,
%  « thợ đóng thùng » l'ouvrier qui monte les tonneaux, « gốc nho » le
%  pied de raisin — et « nho », le raisin, est lui-meme un vieil
%  emprunt au chinois. C'est l'inverse exact du tableau 7, ou la cour
%  de ferme etait chez elle.
%
%  « BLEUET » NE POSE AUCUN PROBLEME ICI, la ou il etait le piege du
%  quebecois : la fleur des champs se dit « hoa thanh cúc » et la baie
%  n'existe pas dans le paysage. Mais la digitale, la bruyere et
%  l'aubepine n'ont pas de nom populaire, et ce sont leurs noms
%  savants qui servent.
% ===================================================================

%%K t08-noto-1 noto
\VUnotes{143.2mm}{%
\textit{(*) Vì từ này không chính xác : đó là mùa thu hoạch lin, mùa hái nho hay mùa hái
táo? Có lẽ nên dùng «} \VUgras{mesono} \textit{» đã được đề nghị trong Mondo XIV, tr.
242. Nhưng đến nay gốc từ mới ấy chưa được Hàn lâm viện chấp nhận và còn chờ được bàn
theo thể lệ thường lệ của giới ido.}%
}

%%K t08-apar-1 apar
\VUsaut{5.0mm}
\VUtitre{15.0pt}{60}{BẢNG Số 8}
\VUsaut{3.0mm}
\VUcentre{13.2pt}{90}{\textit{Cảnh thứ nhất.}}
\VUsaut{3.0mm}
\VUpk{10.2pt}{40}{Mùa gặt (*).}
\VUsaut{4.0mm}
\VUpk{10.2pt}{40}{Cảnh sắc đồng quê.}
\VUsaut{3.0mm}

%%K t08-c1-01-1 p
1. --- Cùng với \VUgras{mùa hè}, cảnh sắc đồng quê lại thay đổi một lần nữa. Hạt giống
gửi vào luống cày đã nảy mầm; một cây nhỏ màu xanh đã mọc lên khỏi mặt đất và trổ bông.
Hạt đã kết, rồi chín, và từ giờ chỉ còn việc gặt.

%%K t08-c1-02-1 p
2. --- \VUgras{Hoa đồng nội} đã nở. \VUgras{Hoa thanh cúc} \textsuperscript{(1)},
\VUgras{hoa anh túc} \textsuperscript{(2)} và \VUgras{hoa mao địa hoàng}
\textsuperscript{(3)} điểm vào màu đơn điệu của những cánh đồng lúa mì nét tương phản
vui tươi của những \VUgras{tràng hoa} tươi mới.

%%K t08-c1-03-1 p
3. --- Cây ăn quả đã có lá; quả đã chín. Cây \VUgras{anh đào} \textsuperscript{(4)} nhỏ
trĩu những \VUgras{quả anh đào} \textsuperscript{(5)} đẹp mà trẻ con thích thú hái và
ăn.

%%K t08-c1-04-1 p
4. --- Từ nay gia súc có thể ở ngoài trời suốt ngày.

%%K t08-c1-04-2 p
Những con \VUgras{cừu con} \textsuperscript{(6)} đã lớn và thành những con \VUgras{cừu
cái} \textsuperscript{(7)} đẹp và những con \VUgras{cừu đực} \textsuperscript{(8)} đẹp
với bộ lông dày. Chúng bình thản gặm \VUgras{cỏ} của \VUgras{bãi chăn}
\textsuperscript{(9)} lấm tấm \VUgras{hoa cúc dại}.

%%K t08-c1-04-3 p
Sắp tới người ta sẽ xén lông \VUgras{những} con vật ấy để lấy \VUgras{len} dệt dạ may
quần áo cho ta.

%%K t08-tit-2 sub
\VUsaut{5.0mm}
\VUpk{10.2pt}{40}{Người chăn cừu.}
\VUsaut{3.4mm}

%%K t08-c1-05-1 p
5. --- \VUgras{Người chăn cừu} \textsuperscript{(10)} có vẻ không mấy để ý đến chúng.
Ông chỉ lo về cơn giông đang đi qua ở chân trời, còn \VUgras{người giữ đàn}
\textsuperscript{(13)}, đang đưa \VUgras{bầu nước} \textsuperscript{(14)} lên môi, thì
có vẻ còn vô tư hơn.

%%K t08-c1-06-1 p
6. --- Trời nóng ngột ngạt, vì con \VUgras{chó} \textsuperscript{(15)} cũng đến giải
khát; nó liếm nước mát của \VUgras{con suối} \textsuperscript{(16)} và làm một con
\VUgras{ếch} \textsuperscript{(17)} nhỏ đáng thương nấp trong đám cỏ ven bờ hoảng sợ. Con
ếch nhảy xuống làn nước trong nơi \VUgras{hoa súng} \textsuperscript{(18)} đang nở.

%%K t08-c1-07-1 p
7. --- Một con \VUgras{chim chìa vôi} \textsuperscript{(19)} tắm gần \VUgras{con mạch
nước} \textsuperscript{(20)} phun ra giữa hai tảng đá và nấp dưới \VUgras{bụi gai}
\textsuperscript{(21)} và \VUgras{bụi táo gai} \textsuperscript{(22)}.

%%K t08-tit-3 sub
\VUsaut{5.0mm}
\VUpk{10.2pt}{40}{Mùa gặt.}
\VUsaut{3.4mm}

%%K t08-c1-08-1 p
8. --- Với trẻ con nhà quê, mùa gặt lúa mì là một thời kỳ rất vui. Chúng vui vẻ
\VUgras{nhào lộn} \textsuperscript{(24)} trên những \VUgras{đống rơm}
\textsuperscript{(23)}, chúng giúp \VUgras{thợ gặt} \textsuperscript{(26)}, và trong khi
những người này cắt \VUgras{ruộng lúa mì} \textsuperscript{(27)} bằng \VUgras{hái}
\textsuperscript{(25)} đã mài sắc trên \VUgras{đá mài} \textsuperscript{(30)}, chúng gom
lúa lại và buộc thành \VUgras{bó} \textsuperscript{(31)} hay thành \VUgras{lượm}
\textsuperscript{(32)}.

%%K t08-c1-09-1 p
9. --- Đôi khi người ta cho những đứa lớn nhất được cắt bằng \VUgras{liềm}
\textsuperscript{(12)} những \VUgras{thân lúa vàng} \textsuperscript{(28)} mang những
\VUgras{bông} \textsuperscript{(29)} nặng trĩu hạt; còn những đứa nhỏ, trong lúc trông
\VUgras{đàn gia súc} \textsuperscript{(33)} đang bình thản gặm cỏ trên \VUgras{bãi chăn}
\textsuperscript{(34)}, dưới bóng cây \VUgras{đoạn} \textsuperscript{(35)} lớn, thì bắt
những con \VUgras{bướm} đẹp bằng \VUgras{vợt} \textsuperscript{(36)} của mình.

%%K t08-c1-09-2 p
Vài đứa còn trèo lên \VUgras{xe} \textsuperscript{(37)} để xếp những bó lúa mà người ta
chuyền lên cho chúng bằng đầu \VUgras{chĩa} \textsuperscript{(38)}.

%%K t08-c1-10-1 p
10. --- Trên khắp \VUgras{cánh đồng} \textsuperscript{(39)}, nơi những con \VUgras{chim
sơn ca} \textsuperscript{(41)} bay lên, không khí thật nhộn nhịp. Ai nấy đều bận việc
gặt. Nhưng trong khi những người nghèo còn dùng các dụng cụ cũ — \VUgras{hái},
\VUgras{liềm} và \VUgras{néo đập} \textsuperscript{(40)} — thì các điền chủ giàu cho cắt
lúa mì hay \VUgras{lúa mạch đen} \textsuperscript{(43)} của mình bằng \VUgras{máy gặt}
\textsuperscript{(42)}. Sau đó, vì không phải chở các bó lúa về kho, người ta chất ngay
tại chỗ những \VUgras{đống lúa} \textsuperscript{(44)} lớn.

%%K t08-c1-11-1 p
11. --- Rồi người ta đưa đến một chiếc \VUgras{máy đập} \textsuperscript{(47)} do một
\VUgras{đầu máy di động} \textsuperscript{(45)} làm chạy nhờ một \VUgras{dây truyền
động} \textsuperscript{(46)}, và như thế hạt được tách khỏi \VUgras{rơm} ngay trên thửa
ruộng nơi nó đã được gieo.

%%K t08-tit-4 sub
\VUsaut{5.0mm}
\VUpk{10.2pt}{40}{Cơn giông.}
\VUsaut{3.4mm}

%%K t08-c1-12-1 p
12. --- Vào mùa gặt thường có những \VUgras{cơn giông} \textsuperscript{(53)} dữ dội.
Trước hết người ta nghe tiếng \VUgras{sấm} ầm ì từ xa; một \VUgras{trận gió tây}
\textsuperscript{(54)} mạnh nổi lên và hổn hển bẻ cong những cây lớn nhất của
\VUgras{khu rừng nhỏ} \textsuperscript{(49)}. Những \VUgras{đám mây}
\textsuperscript{(56)} đen ùn ùn kéo tới; chúng bị những đường ngoằn ngoèo chói lòa của
\VUgras{tia chớp} \textsuperscript{(55)} xuyên qua. \VUgras{Mưa} và đôi khi cả
\VUgras{mưa đá} bắt đầu rơi, đè bẹp những vụ mùa sắp được gặt.

%%K t08-c1-13-1 p
13. --- Sét thường làm cháy một công trình. Chẳng hạn năm ngoái, mái \VUgras{cối xay
gió} \textsuperscript{(50)} đã bị thiêu ra tro và hai chiếc \VUgras{cánh}
\textsuperscript{(51)} của nó bị gãy nát.

%%K t08-c1-14-1 p
14. --- Sau vài giờ, sự yên tĩnh trở lại. Một chiếc \VUgras{cầu vồng}
\textsuperscript{(52)} rực rỡ hiện ở chân trời và thiên nhiên lấy lại nhịp sống đã bị
ngắt quãng trong chốc lát. Những người nhà quê đã trú trong các \VUgras{túp lều}
\textsuperscript{(48)} trở lại làm việc và can đảm cố gắng sửa chữa những thiệt hại vừa
phải chịu.

%%K t08-tit-5 sub
\VUsaut{6.0mm}
\VUcentre{13.2pt}{90}{\textit{Cảnh thứ hai.}}
\VUsaut{3.6mm}

%%K t08-tit-6 sub
\VUpk{10.2pt}{40}{Đi săn. --- Mùa hái nho.}
\VUsaut{1.4mm}
\VUpk{10.2pt}{40}{Đi câu.}
\VUsaut{4.0mm}

%%K t08-c2-01-1 p
1. --- Vào cuối \VUgras{tháng tám} hay giữa \VUgras{tháng chín}, tùy vùng, là ngày mở
đầu mùa săn. Những tia nắng đầu tiên vừa làm những con \VUgras{thằn lằn}
\textsuperscript{(2)} ra khỏi hang thì \VUgras{người đi săn} \textsuperscript{(5)} đã
lên đường cùng đàn \VUgras{chó} \textsuperscript{(4)} của mình.

%%K t08-c2-02-1 p
2. --- Ông mặc bộ \VUgras{đồ đi săn} \textsuperscript{(9)} chắc chắn bằng dạ dày, đi
đôi \VUgras{ghệt} da giúp ông lội được trong đám cỏ ướt nơi những con \VUgras{cóc}
\textsuperscript{(1)} ẩn mình; ông cầm khẩu \VUgras{súng} \textsuperscript{(6)} và chiếc
\VUgras{túi săn} \textsuperscript{(10)}, và thế là ông đi.

%%K t08-c2-03-1 p
3. --- Cuộc đuổi bắt hăng say đưa ông đến giữa một vườn nho nơi mùa hái nho đang rộn
ràng. Con cái người trồng nho, thường ngày vẫn giữ dê bên đường, hôm nay để chúng ăn cỏ
tự do và, sau khi buộc vào một cái cọc con \VUgras{dê đực} \textsuperscript{(3)} già
chắc chắn sẽ lợi dụng tự do mà chạy mất, chúng cũng tự dành cho mình một phần thú vui.
Một đứa lấy một chùm nho trong \VUgras{giỏ hái nho} \textsuperscript{(12)} và thích thú
vắt \VUgras{nước} \textsuperscript{(14)} vào một chiếc \VUgras{cốc}
\textsuperscript{(13)} để làm nước nho tươi (rượu chưa lên men). Đứa kia đội lên đầu một
\VUgras{vòng lá} \textsuperscript{(15)} vừa tết xong. Bên cạnh chúng, những con \VUgras{chim
sẻ} \textsuperscript{(16)} dạn dĩ đến tận trước máy ép để mổ trộm nho.

%%K t08-c2-04-1 p
4. --- Trong lúc trẻ con chơi đùa, người lớn làm việc. \VUgras{Thợ hái nho}
\textsuperscript{(18)} mang nho về hầm bằng những chiếc \VUgras{gùi}
\textsuperscript{(17)}; một \VUgras{thợ đóng thùng} \textsuperscript{(19)} sửa các
\VUgras{thùng rượu} \textsuperscript{(20)}, kiểm tra xem \VUgras{ván thùng}
\textsuperscript{(22)} và \VUgras{đáy thùng} \textsuperscript{(21)} có tốt không, và
thay những chiếc \VUgras{đai} \textsuperscript{(23)} gãy. Chính ông cũng đã làm những
\VUgras{bồn lớn} \textsuperscript{(25)} mà người ta đổ \VUgras{nho}
\textsuperscript{(26)} vào cho lên men.

%%K t08-c2-05-1 p
5. --- Khi lên men xong, người ta rút rượu ra và đặt bã lên \VUgras{máy ép}
\textsuperscript{(28)}, rồi vặn dần chiếc \VUgras{trục vít} \textsuperscript{(29)} để ép
lấy nước, nước chảy vào một cái \VUgras{chậu} \textsuperscript{(32)}, và người ta lại
được thêm \VUgras{rượu vang} \textsuperscript{(31)}.

%%K t08-tit-8 sub
\VUsaut{6.0mm}
\VUpk{10.2pt}{40}{Bia và rượu táo.}
\VUsaut{3.4mm}

%%K t08-c2-06-1 p
6. --- Trên tường \VUgras{hầm rượu} \textsuperscript{(27)} leo một nhánh \VUgras{hoa
bia} \textsuperscript{(30)}. Ở đây nó chỉ là cây cảnh, nhưng ở vài nước, nơi cây nho rất
hiếm, người ta trồng hoa bia để làm \VUgras{bia}.

%%K t08-c2-07-1 p
7. --- Cây \VUgras{táo} \textsuperscript{(33)} nhỏ trĩu quả, những quả ấy khi chín sẽ cho
thứ \VUgras{rượu táo} tuyệt hảo. Trong \VUgras{lồng} \textsuperscript{(34)} của chúng,
con \VUgras{chim hoàng yến} \textsuperscript{(35)} và con \VUgras{chim kim oanh}
\textsuperscript{(36)} nhìn bằng con mắt thèm thuồng những con chim sẻ đang bay nhảy tự
do.

%%K t08-c2-08-1 p
8. --- Ngút tầm mắt, trên sườn đồi, thẳng hàng những \VUgras{cọc chống}
\textsuperscript{(40)} đỡ các \VUgras{gốc nho} \textsuperscript{(39)} tuyệt đẹp, trĩu
những \VUgras{chùm} nặng.

%%K t08-c2-08-2 p
Suốt cả năm, \VUgras{người trồng nho} \textsuperscript{(42)} đã làm lụng để chăm sóc
\VUgras{vườn nho} \textsuperscript{(38)} của mình, và giờ đây ông được đền công bằng một
vụ mùa đẹp. Những \VUgras{cô hái nho} \textsuperscript{(41)} đổ đầy các bồn đặt trên
chiếc xe do đàn \VUgras{la} \textsuperscript{(43)} kéo, và ai nấy đều vui.

%%K t08-tit-9 sub
\VUsaut{5.0mm}
\VUpk{10.2pt}{40}{Đi câu.}
\VUsaut{3.4mm}

%%K t08-c2-09-1 p
9. --- Những \VUgras{người câu cá} \textsuperscript{(44)} cũng vậy; từ sáng sớm họ đã
đến ngồi bên bờ nước với những chiếc \VUgras{cần câu} \textsuperscript{(45)} của mình.

%%K t08-c2-09-2 p
\VUgras{Cá} \textsuperscript{(47)} cắn \VUgras{lưỡi câu} ngay khi họ thả \VUgras{dây
câu} \textsuperscript{(46)} xuống nước, và họ vừa kịp thay \VUgras{mồi} thì một con nữa
đã giãy giụa ở đầu dây.

%%K t08-c2-09-3 p
Tuy vậy, \VUgras{người quăng chài} \textsuperscript{(48)}, đứng trên chiếc \VUgras{thuyền}
\textsuperscript{(49)} của mình quăng \VUgras{chài} ra giữa \VUgras{dòng sông}
\textsuperscript{(50)}, bắt được nhiều cá hơn.

%%K t08-c2-10-1 p
10. --- Mùa thu là mùa của những chuyến dạo chơi đẹp : đi bộ, đi \VUgras{ngựa}
\textsuperscript{(52)}, đi \VUgras{xe đạp} \textsuperscript{(53)}, đi \VUgras{ô-tô}
\textsuperscript{(54)}. \VUgras{Đường sá} \textsuperscript{(51)} sạch sẽ và người ta tận
hưởng những ngày đẹp cuối cùng, vì kìa những con \VUgras{chim én}
\textsuperscript{(56)} đã tụ trên mái nhà để bay đi hết cánh về những xứ ấm hơn. Lá của
cây \VUgras{hồ đào} \textsuperscript{(57)} già đang ngả vàng, những quả \VUgras{hồ đào}
rụng xuống, và những \VUgras{cây} cao mọc thành \VUgras{cụm} \textsuperscript{(58)} trên
\VUgras{gò cao} \textsuperscript{(59)} sắp trơ cành.

%%K t08-c2-11-1 p
11. --- \VUgras{Mặt trời} \textsuperscript{(60)} mọc muộn hơn, ngày ngắn lại, một cái
lạnh khá buốt bắt đầu thấy được vào \VUgras{buổi sáng}, và kìa những đàn \VUgras{gà rừng}
\textsuperscript{(61)} đã rời xứ ta, nay đã kém hiếu khách.
\VUsaut{6.0mm}
\VUfilet{22mm}
